\documentclass[10pt]{article}
%% Shared preamble for the 9.520 Oct 2 literature reviews.
\usepackage[T1]{fontenc}
\usepackage[utf8]{inputenc}
\usepackage{lmodern}
\usepackage[margin=0.75in]{geometry}
\usepackage{amsmath,amssymb}
\usepackage{microtype}
\usepackage{enumitem}
\usepackage[dvipsnames]{xcolor}
\usepackage[backend=bibtex,style=authoryear,natbib=true,maxcitenames=2,maxbibnames=4,minbibnames=3,uniquelist=false,uniquename=false,giveninits=true,doi=false,isbn=false,url=false,eprint=true]{biblatex}
\addbibresource{refs.bib}
\usepackage[colorlinks=true,linkcolor=MidnightBlue,citecolor=MidnightBlue,urlcolor=MidnightBlue]{hyperref}
\usepackage[capitalise,nameinlink]{cleveref}
\setlength{\parskip}{2pt}
\setlength{\parindent}{0pt}
\setlist{nosep,leftmargin=*}

%% Tags for the annotated entries: [F] foundational, [R] 2024-26, [P] preprint or workshop.
\newcommand{\tagF}{\textsf{\textcolor{MidnightBlue}{[F]}}}
\newcommand{\tagR}{\textsf{\textcolor{ForestGreen}{[R]}}}
\newcommand{\tagP}{\textsf{\textcolor{BrickRed}{[P]}}}

%% One annotated bibliography entry: \entry{key}{tags}{annotation}
\DeclareBibliographyCategory{annotated}
\newcommand{\entry}[3]{%
  \addtocategory{annotated}{#1}%
  \par\noindent\hangindent=1.2em\hangafter=1
  {\small\fullcite{#1}}~#2\par
  {\leftskip=1.2em #3\par}\vspace{2pt}}

%% Group synthesis paragraph opening each explanation.
\newcommand{\synth}[1]{\par\noindent\emph{#1}\par\vspace{2pt}}

%% Works cited in the text but not annotated, printed compactly at the end.
\newcommand{\otherrefs}{%
  {\footnotesize\setlength{\bibitemsep}{0pt}%
  \printbibliography[notcategory=annotated,heading=subbibliography,title={Other references}]}}


\title{\vspace{-1.6em}Why Does Late-Stage Learning-Rate Annealing Help?\\[2pt]
\large Literature review and implications (9.520, Fall 2026)}
\author{Juntang Wang \and V\'ictor Conchello}
\date{October 2, 2026}

\begin{document}
\maketitle
\vspace{-1.5em}

%% =================================================================
%% Question
%% =================================================================
\section{The question and why it is hard}
\label{sec:question}

Most large runs end with a cooldown: the loss plateaus at a constant learning rate $\eta$, then drops
sharply once $\eta$ decays, far more than the same steps at constant $\eta$ would give
\citep{hagele2024scaling,wen2025river}. Four explanations compete. The decay may (i) restore a
\emph{stability margin}, pulling the iterate off the edge of stability (EoS); (ii) \emph{reduce
gradient noise}; (iii) \emph{select a basin} the constant run never reaches; or (iv) just add
\emph{compute} at a useful moment. One decay changes all four at once, and a loss curve fits each
story. Telling them apart needs interventions that move one factor at a matched update budget. Each
group below collects the evidence for one explanation. Tags: \tagF{} foundational, \tagR{} 2024--26
top venue, \tagP{} preprint or workshop.

%% =================================================================
%% Literature
%% =================================================================
\section{Annotated literature}
\label{sec:lit}

\subsection{Stability margin}
\label{sec:lit-stability}

\synth{The only curvature-based account. The first four papers fix $\eta$, so applying them to
annealing is our extrapolation: at EoS the observed loss includes an oscillation term set by $\eta$,
and a decay should remove it. Only the fifth studies the decay.}

\entry{cohen2021eos}{\tagF}{In full-batch GD, the top Hessian eigenvalue $\lambda_{\max}$ rises
(\emph{progressive sharpening}) until it reaches $2/\eta$, then hovers there while the loss keeps
falling non-monotonically, against the classical requirement $\eta<2/\lambda_{\max}$. For annealing:
after a drop to $\eta'$, sharpness may climb again to $2/\eta'$, so the extra margin is temporary.}

\entry{damian2023selfstab}{\tagF}{A cubic Taylor term explains EoS: when the iterate diverges along
the top eigenvector, the cubic term lowers curvature until stability returns. GD at EoS thus follows
projected GD under $S(\theta)\le 2/\eta$. A decay loosens this constraint and lets the trajectory
move to sharper regions.}

\entry{cohen2025central}{\tagR}{Derives a \emph{central flow}, an ODE for the time-averaged
trajectory at EoS that predicts long-run loss and sharpness, also for adaptive optimizers. The
observed loss exceeds the loss at the time-averaged point by an amount set by the oscillations; that
gap is our candidate for the immediate cooldown drop. Schedules are not studied.}

\entry{andreyev2024eoss}{\tagP}{For mini-batch SGD, what settles at $2/\eta$ is \emph{batch
sharpness}, the curvature of mini-batch losses along their own gradients, while $\lambda_{\max}$,
usually smaller, is suppressed. So with SGD it is batch sharpness that should re-equilibrate after a
decay.}

\entry{wen2025river}{\tagR}{Pretraining loss as a \emph{river valley}: at high $\eta$ the iterate
moves fast along the river while bouncing between the walls; the decay stops the bouncing and reveals
the progress. Both phases are needed. The landscape is a conjecture, shown to arise on a bigram
dataset.}

\subsection{Noise reduction}
\label{sec:lit-noise}

\synth{A constant step size leaves an error floor proportional to $\eta$, and a decay removes it.
Caveat: in non-smooth convex theory the same floor comes from gradient size even without noise, so
convex bounds fitting LLM curves do not by themselves favour noise over oscillation.}

\entry{smith2018dontdecay}{\tagF}{Raising the batch size instead of decaying $\eta$ gives the same
training and test curves, for SGD, momentum and Adam, with fewer updates. If annealing were only
noise reduction, this swap would be exact. Shown for step decay in image classification; it also
changes the update count.}

\entry{ge2019stepdecay}{\tagF}{For the final SGD iterate on least squares with a known horizon, every
polynomial decay is far from the minimax rate, while geometric step decay is within log factors. The
shape of the decay matters even when noise is the only reason to decay.}

\entry{schaipp2025convex}{\tagR}{A last-iterate bound from non-smooth convex optimization tracks LLM
loss curves: with a linear cooldown it loses its log term, matching the observed gain, and it
transfers good learning rates across schedules. Its floor comes from gradient norms, so it fits (ii)
and an oscillation reading of (i) equally well.}

\entry{li2026fsl}{\tagR}{Derives the optimal schedule under a stability constraint for SGD in
power-law kernel regression. Easy tasks: decay from the start. Hard tasks: hold the largest stable
$\eta$, then decay at the end, i.e.\ WSD. Difficulty sets when to decay; capacity sets how.}

\subsection{Basin selection}
\label{sec:lit-basin}

\synth{Here the high-$\eta$ phase regularizes and the decay reveals the result. Two papers concern
test error, not training loss. One recent paper finds the opposite sign: the basin after a decay can
be worse for later use.}

\entry{li2019largelr}{\tagF}{Proves in a two-layer model that large initial $\eta$ plus annealing
beats small $\eta$ throughout. The order of learning decides it: the small-$\eta$ model first
memorizes easy-to-generalize, hard-to-fit patterns, then generalizes worse on hard-to-generalize,
easier-to-fit ones. On CIFAR-10, a small patch is memorized at once at small $\eta$ but ignored at
large $\eta$ until after annealing.}

\entry{andriushchenko2023sparse}{\tagF}{Large steps make SGD bounce across a valley, which stalls the
loss; meanwhile noise orthogonal to the bouncing pushes the network toward sparse predictors. The
longer $\eta$ stays high, the stronger the effect, so the plateau does real work that the decay
exposes. Same geometry as \citet{wen2025river}, read as regularization.}

\entry{kwag2026eos}{\tagP}{From the instructors' group. Branching runs that enter or leave EoS from
the same state show that the stability constraint shifts learning across data subsets: groups whose
gradient aligns with the top Hessian eigenvector, and does not saturate, gain at others' expense. A
decay is one such exit, so it should change \emph{which} examples get fit, which the average loss
hides. Their branching design is a template for our ablation.}

\entry{yano2026nodecay}{\tagR}{At 1B and 8B parameters, skipping the final decay can give worse
pretraining loss but better results after supervised fine-tuning. Decay lands in sharper minima.
Whether annealing ``helps'' depends on the metric, and the sharper minima fit the stability view.}

\subsection{Extra compute}
\label{sec:lit-compute}

\synth{These compare schedules at matched budgets, the closest existing controls for (iv). Three find
the cooldown gain survives at equal compute; two find that averaging the iterates recovers much of
it, which points to noise or oscillation, not extra steps. Iterate averaging goes back to SWA
\citep{izmailov2018swa}.}

\entry{hagele2024scaling}{\tagR}{Constant $\eta$ plus a cooldown scales as predictably as cosine and
matches it at equal compute; weight averaging improves iterates at no training cost. One constant run
can be branched into many cooldowns, so equal-step comparisons are cheap.}

\entry{tissue2025annealing}{\tagR}{Fits LLM loss curves as $L(s)=L_0+A\,S_1^{-\alpha}-C\,S_2$, with
$S_1$ the area under the learning-rate curve and $S_2$ an annealing area, and predicts loss under new
schedules from one or two runs. $S_2$ splits annealing from compute, but descriptively.}

\entry{defazio2024schedulefree}{\tagR}{Schedule-free SGD and AdamW, built on a theory linking
schedules to iterate averaging, need no decay or stopping time and match tuned schedules from convex
problems to large models. If averaging can replace the decay, the gain is mostly variance or
oscillation removal, not a new basin.}

\entry{bergsma2025straight}{\tagR}{At compute-optimal budgets, linear decay to zero beats cosine to
10\% of peak. Reading AdamW's weights as a moving average of updates, the decay widens the window
and averages out noise. The gain grows with tokens per parameter, which predicts where it fails:
short runs, with little noise to average.}

%% =================================================================
%% Key papers
%% =================================================================
\section{Three key papers}
\label{sec:key}

\textbf{\citet{cohen2021eos}} (ICLR; about 500 citations on Semantic Scholar) set up the
stability-margin explanation and the course's training-dynamics line. It replaced
$\eta<2/\lambda_{\max}$ with an empirical law, sharpening then self-stabilization at $2/\eta$, which
self-stabilization, central flows and EoSS all try to explain. Its prediction for us is sharp: after a
decay, sharpness should rise to the new threshold.

\textbf{\citet{wen2025river}} (ICLR; about 80 citations) is one of few top-venue papers with a
mechanism for the decay phase itself. The river valley links the WSD literature to EoS, since the
oscillation across the valley is what the decay removes, and \citet{andriushchenko2023sparse}
describe the same geometry independently. It is falsifiable: progress along the river should
continue while the loss looks flat.

\textbf{\citet{hagele2024scaling}} (NeurIPS spotlight; about 160 citations) gives the matched-compute
controls the project asks for. Branching cooldowns off one run compares schedules at equal steps,
shows the decay carries the gain regardless of the cosine shape, and shows averaging recovers part of
it. The protocol is now common in scaling studies and is the cheapest design for our ablation.

\clearpage
%% =================================================================
%% Implications
%% =================================================================
\section{Implications for theory and practice}
\label{sec:implications}

\begin{enumerate}[leftmargin=1.3em,itemsep=4pt]
\item \textbf{Existing interventions already separate the four explanations.} Full-batch GD has no
noise, so its drop must come from stability or basin selection; a constant-$\eta$ branch of equal
length controls compute (\cref{sec:lit-stability,sec:lit-basin}). Raising $B$ at fixed $\eta$
\citep{smith2018dontdecay} moves only noise. Averaging constant-$\eta$ iterates
\citep{hagele2024scaling,defazio2024schedulefree} removes oscillation and noise without a new basin.
Basin selection needs a basin-sensitive test, not a loss curve: re-warming restores the oscillation
term, so the loss can return to the plateau from a different basin. Instead, decay, re-warm and decay
again, then compare with the directly decayed branch by the loss barrier along the linear path
between them \citep{frankle2020lmc} and by prediction agreement. Run on two tasks at a matched update budget, this grid is the ablation the project asks for,
and the papers above predict every cell.

\item \textbf{Theory: annealing should land in a sharper region, at a cost.} Self-stabilization and
EoSS suggest that after a drop to $\eta'$, sharpness (batch sharpness for SGD) climbs toward
$2/\eta'$ \citep{damian2023selfstab,andreyev2024eoss}. The decayed model then sits somewhere sharper
than the plateau model, matching the fine-tuning penalty in \citet{yano2026nodecay}. ``Annealing
finds a better minimum'' becomes a signed prediction: lower training loss, higher curvature, worse
adaptability. If sharpness does not rise after the decay while the loss keeps improving, the
stability account is wrong for that regime.

\item \textbf{Practice: each substitute for annealing covers only its own component.} A larger batch
replaces the noise component but not the oscillation, because batch sharpness still settles at
$2/\eta$ \citep{andreyev2024eoss}. Averaging removes both but cannot change the basin. So ``increase
the batch'' should fail when the gain is mostly oscillation, as in near full-batch training, and
averaging should fail when the decay changes which features are learned
\citep{andriushchenko2023sparse}. Measuring the gradient noise scale and the oscillation amplitude
before the cooldown would say which tool to use.

\item \textbf{Open question: EoS with a time-varying step size.} Every EoS result we found fixes
$\eta$. A central flow with $\eta(t)$ would predict the whole cooldown: the immediate loss of the
oscillation term \citep{cohen2025central}, then slower re-sharpening. With batch sharpness it would
also bound how fast a schedule may decay before the iterate leaves the stable set. That would give a
first-principles version of the schedule rules \citet{schaipp2025convex} and \citet{li2026fsl} derive
from convex and kernel models, and would test whether the WSD shape comes from the landscape or from
the noise.
\end{enumerate}

\otherrefs

\vspace{2pt}
{\footnotesize\textbf{Use of generative AI.} We used an AI assistant for literature discovery,
checking citation metadata against arXiv and proceedings pages, and drafting prose. We read the
primary sources for every claim above, verified all citations, and take responsibility for the
content.}

\end{document}
